\chapter{Vendoring: Pulling In Other People's Cores}
\label{ch:vendoring}

\chapabstract{Where the code behind a \kw{depend} entry in a \file{.core}
file actually comes from: a pinned, reproducible snapshot copied in by a
small Python tool, not a git submodule.}

\section{What ``vendoring'' means here}

A \file{.core} file can declare \kw{depend: pulp-platform.org::obi}, but
that only works once the files of that dependency actually exist somewhere
\FuseSoC can find them. \textbf{Vendoring} is how they get there: a script
fetches a pinned snapshot of an upstream repository and copies it, as plain
files, into a \file{vendor/} directory inside the lab. No git submodule, no
\file{.gitmodules}, and after the copy, no connection back to the upstream
repository's own \file{.git} history at all.

\begin{ruleofthumb}
Vendoring answers ``where do these files come from and which exact version
are they''. \FuseSoC's \kw{depend} answers ``which files does this target
need''. The two are independent: \FuseSoC does not fetch anything, and the
vendoring script does not know what a target is.
\end{ruleofthumb}

\begin{notebox}{Why not a submodule?}
A submodule keeps the upstream's whole history and its own \file{.git},
which a student clone does not need and a build tool should not have to
walk. A vendored copy is just files: reproducible from the pinned
\kw{rev}, easy to diff, and impossible to accidentally \code{git push}
into.
\end{notebox}

\section{Anatomy of a .vendor.hjson}

Every dependency the lab vendors has its own descriptor, ending in
\file{.vendor.hjson}. Two real ones, at opposite ends of the size scale.

\begin{corecode}[caption={lab1/vendor/pulp\_platform/common\_cells.vendor.hjson},label={lst:vendoring:common}]
// Vendoring descriptor: pulp-platform common_cells.
// Fetched by `make vendor` (python3 util/vendor.py <this file>).
{
  name: "common_cells",
  target_dir: "common_cells",

  upstream: {
    url: "https://github.com/pulp-platform/common_cells.git",
    rev: "v1.38.0",
  },

  exclude_from_upstream: [
    "test",
    "ci",
    "formal",
    ".github",
  ]
}
\end{corecode}

\begin{corecode}[caption={lab0/x-heep.vendor.hjson (the whole file)},label={lst:vendoring:xheep}]
// Vendoring descriptor for X-HEEP.
// `make vendor` runs: python3 util/vendor.py x-heep.vendor.hjson
{
  name: "x-heep",
  target_dir: "x-heep",

  upstream: {
    url: "https://github.com/esl-epfl/x-heep.git",
    rev: "main",
  },
}
\end{corecode}

The fields are the same regardless of size:

\begin{itemize}
  \item \kw{name}: a label for the descriptor itself, not a \FuseSoC core
        name.
  \item \kw{target\_dir}: where the snapshot lands, relative to
        \file{vendor/} (or wherever the enclosing Makefile passes it).
  \item \kw{upstream}: \kw{url} plus \kw{rev} -- a tag, branch, or commit
        SHA. This is the whole reproducibility story: two people running
        \code{make vendor} against the same descriptor get byte-identical
        sources.
  \item \kw{exclude\_from\_upstream} (optional): paths to drop from the
        copy -- typically the upstream's own tests and CI, which the lab
        does not need and does not want cluttering \file{vendor/}.
\end{itemize}

\begin{gotcha}{\code{rev: "main"} is a moving target}
\Cref{lst:vendoring:xheep} pins to \code{main} on purpose while the lab is
still being developed: it always vendors the latest X-HEEP. Once a lab is
handed to students, \code{rev} should be pinned to a commit SHA -- otherwise
two students running \code{make vendor} on different days can end up with
different X-HEEP versions and different bugs.
\end{gotcha}

\section{What make vendor actually does}

\begin{shcode}
python3 util/vendor.py vendor/pulp_platform/common_cells.vendor.hjson
\end{shcode}

The script (\file{util/vendor.py}, itself vendored from OpenTitan and not
something you need to read) does three things: fetches the repository at the
pinned \kw{rev} into a temporary location, copies everything except
\kw{exclude\_from\_upstream} into \kw{target\_dir}, and writes a
\file{<name>.lock.hjson} next to the descriptor recording exactly what was
fetched (the resolved commit, the date). The lock file is what you would
diff to answer ``did this dependency change since last time''.

A lab's Makefile runs this once per descriptor, over every
\file{*.vendor.hjson} it finds:

\begin{makecode}[caption={lab1/Makefile -- the vendor target}]
VENDOR_HJSON := $(wildcard $(VENDOR)/*.vendor.hjson)

.PHONY: vendor
vendor:
	@for f in $(VENDOR_HJSON); do \
	  echo "$(PYTHON) util/vendor.py $$f"; \
	  $(PYTHON) util/vendor.py $$f || exit 1; \
	done
\end{makecode}

\section{Two vendoring patterns in these labs}

Lab 1 vendors several small IPs individually, each with its own
\file{.vendor.hjson} under \file{vendor/pulp\_platform/}, because
\file{cordic\_accel.core} depends on each of them by name separately. Lab 0
vendors one thing, X-HEEP itself, as a single large snapshot, because the
whole lab is built on top of it rather than depending on named pieces of
it. Neither is more correct: vendor at the granularity your \file{.core}
files actually depend on.

\section{Checklist}
\label{sec:vendoring:checklist}

\begin{itemize}
  \item Vendoring copies pinned files in; it is not a submodule and keeps no
        upstream history.
  \item \kw{upstream.rev} is the whole reproducibility guarantee -- pin it
        to a SHA or tag before handing a lab out.
  \item \file{util/vendor.py <descriptor>.vendor.hjson} does the fetch;
        \code{make vendor} just runs it over every descriptor it finds.
  \item The generated \file{.lock.hjson} records what was actually fetched;
        do not hand-edit it.
  \item A vendored directory is regenerated by \code{make vendor}, so it is
        normal for it to be absent from a fresh clone and for
        \code{.gitignore} to exclude it.
\end{itemize}
